\chapter{EEG}
\index{EEG}

\section*{Definition and Overview}
An electroencephalogram, known as an EEG, is a test that records the electrical activity of the brain. Electrodes placed on the scalp detect the tiny electrical signals produced by brain cells, which are displayed as wavy lines on a computer screen. The EEG is used mainly to investigate epilepsy and seizures, but it also helps in the assessment of sleep disorders, confusion, brain injury, and other neurological conditions. It is a safe, painless test with no radiation and no electrical currents passed into the head.

\section*{Epidemiology}
EEG testing is performed widely in many countries for people with suspected or known epilepsy, and it is an important tool in hospital neurology and sleep medicine. Epilepsy affects about 1 in 100 people, and most of those with new seizures will have an EEG as part of their assessment. The test is also used in the intensive care unit, in sleep laboratories, and in the assessment of dementia and encephalitis.

\section*{Aetiology and Pathophysiology}
The brain's electrical activity arises from the coordinated firing of large numbers of nerve cells.
\begin{itemize}
    \item \textbf{Brain waves:} the EEG records rhythmic patterns of electrical activity, described by frequency as alpha, beta, theta, and delta waves
    \item \textbf{Normal rhythms:} the dominant rhythm varies with age, wakefulness, and mental state, and changes during sleep
    \item \textbf{Epileptic activity:} in epilepsy, bursts of abnormal electrical discharge produce spikes and sharp waves on the EEG
    \item \textbf{Other abnormalities:} slowing of brain waves occurs with metabolic disturbance, infection, drugs, and brain injury
    \item \textbf{Activation:} flashing lights, hyperventilation, and sleep can provoke epileptic activity during the recording
    \item \textbf{Limitations:} a normal EEG does not exclude epilepsy, and some people without epilepsy have minor non-specific changes
\end{itemize}

\section*{Clinical Features}
An EEG is arranged to answer specific clinical questions.
\begin{itemize}
    \item Investigation of suspected seizures or epilepsy
    \item Classification of the type of epilepsy to guide treatment
    \item Assessment of unusual episodes that may be seizures, faints, or other events
    \item Investigation of confusion, delirium, or altered consciousness
    \item Assessment of sleep disorders, including sleep-related epilepsy
    \item Monitoring during surgery or in the intensive care unit
    \item Brain death testing in specific circumstances
\end{itemize}

\section*{Differential Diagnosis}
The EEG is a diagnostic tool whose interpretation depends on the clinical question.
\begin{itemize}
    \item Seizures may originate in brain regions that do not show changes on a routine scalp EEG
    \item Fainting, panic attacks, and movement disorders can mimic seizures and need to be distinguished clinically
    \item Some people without epilepsy have EEG features that can be misinterpreted
    \item The result is always combined with the history and examination
\end{itemize}

\section*{When to Seek Medical Care}
An EEG is arranged by a specialist after clinical assessment. Anyone who has had a possible seizure, an episode of loss of consciousness, or unexplained confusion should seek medical advice.

\textbf{Contact your local emergency services for:}
\begin{itemize}
    \item A seizure lasting more than five minutes, or repeated seizures without recovery
    \item A first seizure with difficulty breathing, or in pregnancy
    \item Prolonged loss of consciousness
    \item Sudden weakness, facial drooping, or difficulty speaking, which may indicate stroke
\end{itemize}

\section*{Diagnosis and Investigations}
Undergoing an EEG is straightforward.
\begin{itemize}
    \item \textbf{Preparation:} hair should be clean and free of oils and products; sleep deprivation may be arranged beforehand to improve the yield
    \item \textbf{The test:} electrodes are attached to the scalp with paste, and the person rests with eyes open and closed; hyperventilation and flashing lights are used to activate the recording
    \item \textbf{Duration:} a routine EEG takes about 30 to 60 minutes
    \item \textbf{Longer recordings:} sleep-deprived EEGs, 24-hour ambulatory EEGs, and video-EEG monitoring in hospital capture episodes over longer periods
    \item \textbf{Reporting:} a neurologist or clinical neurophysiologist interprets the recording
\end{itemize}

\section*{Management}
The EEG informs, but does not replace, clinical diagnosis and management.
\begin{itemize}
    \item \textbf{Epilepsy diagnosis:} a clear epileptic discharge supports the diagnosis and guides medicine choice
    \item \textbf{Seizure classification:} the pattern helps classify focal versus generalised epilepsy
    \item \textbf{Treatment decisions:} anti-epileptic medicines are chosen partly based on the EEG findings
    \item \textbf{Monitoring:} repeat EEGs may be used to assess treatment or clarify the diagnosis
    \item \textbf{Normal results:} the clinician weighs the EEG against the clinical picture
\end{itemize}

\section*{Complications}
\begin{itemize}
    \item The test is safe, with no significant physical risks
    \item Skin irritation from the electrode paste in a small number of people
    \item Misinterpretation, which is why the EEG is read by specialists in clinical context
    \item The inconvenience of longer recordings for some people
\end{itemize}

\section*{Prognosis and Outcomes}
The EEG is a valuable and reliable test when used and interpreted appropriately. It helps establish the diagnosis of epilepsy, guides treatment, and clarifies the nature of unexplained episodes. Its limitations are well understood by clinicians, and the test contributes to the accurate management of many neurological conditions.

\section*{Prevention and Self-Care}
\begin{itemize}
    \item Wash hair before the test as instructed and avoid hair products
    \item Follow advice about sleep deprivation or medication before the test
    \item Bring a list of medicines and a description of any episodes
    \item Ask what the results mean and how they affect your care
    \item Do not stop epilepsy medicines without medical advice, even before a test
\end{itemize}

\section*{Sources and Further Reading}
The following reputable sources provide further information for healthcare students and patients seeking reliable, current advice about EEG.
\begin{itemize}
    \item Epilepsy Australia. \textit{Information on epilepsy and EEG testing.} https://www.epilepsyaustralia.net/
    \item Healthdirect Australia. \textit{Electroencephalogram (EEG).} https://www.healthdirect.gov.au/eeg
    \item NHS. \textit{Electroencephalogram (EEG).} https://www.nhs.uk/conditions/electroencephalogram/
    \item Mayo Clinic. \textit{EEG (electroencephalogram).} https://www.mayoclinic.org/tests-procedures/eeg/
    \item MedlinePlus. \textit{Electroencephalogram.} https://medlineplus.gov/ency/article/003931.htm
    \item Better Health Channel. \textit{EEG brain function test.} https://www.betterhealth.vic.gov.au/
\end{itemize}
